\subsection{GLR under Graph Inaccuracies}
\label{subsec:problem-framing}

We aim to design a regularizer that promotes graph smoothness using a potentially inaccurate proxy graph while retaining the useful information contained in the available graph structure.
We neither refine the proxy graph itself nor assume prior knowledge of which graph relations are inaccurate, since such information is generally unavailable in practice.
To identify the regularization properties required under such graph uncertainty, we first examine how graph inaccuracies affect GLR through its edge-domain expression.

Although GLR is typically written in the node domain as the quadratic form of the graph Laplacian, it can equivalently be expressed in the edge domain as
\begin{equation}
R_{\mathrm{GLR}}(\mathbf{u})
=
\left\|
\sqrt{\mathbf{W}}\mathbf{D}\mathbf{u}
\right\|_2^2,
\label{eq:glr-edge-domain}
\end{equation}
which follows from $\mathbf{L}=\mathbf{D}^{\top}\mathbf{W}\mathbf{D}$, as stated in Section~II.
Here, $\mathbf{D}$ and $\mathbf{W}$ represent the edge connectivity and edge weights of the available proxy graph, respectively.
Accordingly, we refer to the entries of $\sqrt{\mathbf{W}}\mathbf{D}\mathbf{u}$ as proxy-weighted signal differences.
Topology inaccuracies alter which signal differences are represented through $\mathbf{D}$, whereas edge-weight inaccuracies alter how strongly those differences are scaled through $\mathbf{W}$.
Although these inaccuracies arise from different graph-side mechanisms, their effects on GLR can thus be interpreted through distortions of the proxy-weighted signal differences.

Accordingly, in the remainder of this section, we analyze the robustness of the regularizer to graph inaccuracies primarily through the proxy-weighted signal differences $\sqrt{\mathbf{W}}\mathbf{D}\mathbf{u}$, rather than directly through $\mathbf{D}$ and $\mathbf{W}$.
Under the quadratic evaluation of GLR, these distortions can lead to two distinct failure modes:
\begin{itemize}
    \item \emph{Excessive-penalty failure mode:}
    spurious edges and overestimated weights can associate large edge weights with large signal differences, producing excessively large proxy-weighted signal differences.
    Their quadratic evaluation can then dominate the GLR cost and undesirably suppress legitimate signal variations.

    \item \emph{Information-loss failure mode:}
    missing edges and underestimated weights can associate small or zero edge weights with informative signal differences, producing weakened or missing proxy-weighted signal differences.
    Consequently, competing observation-consistent signals can receive insufficient regularization and become difficult to distinguish from the underlying signal.
\end{itemize}

In this paper, we focus on the excessive-penalty failure mode and consider a magnitude-adaptive smoothness regularization design that mitigates the excessive influence of proxy-weighted signal differences without identifying or correcting inaccurate graph relations themselves.
Although the proposed regularizer does not directly address information loss caused by missing edges or underestimated weights, we also consider mixed-inaccuracy settings in the experiments to assess whether its benefit is retained when both failure modes coexist.


\subsection{Magnitude-Adaptive Regularization via Infimal Convolution}
\label{subsec:magnitude-adaptive}

Building on the preceding analysis, our goal is to design a regularizer that preserves the GLR evaluation as much as possible while preventing excessively large proxy-weighted signal differences from dominating the quadratic evaluation.
To this end, we use infimal convolution to decompose the proxy-weighted signal differences into two components, one evaluated quadratically as in GLR and the other redirected to an alternative penalty when it would otherwise incur an excessively large quadratic cost.

\begin{definition}[Proposed regularizer]
\label{def:proposed-regularizer}
For $\tau>0$, the proposed regularizer is defined as
\begin{equation}
R_{\mathrm{IC}}(\mathbf{u})
=
\Psi_{\tau}
\left(
\sqrt{\mathbf{W}}\mathbf{D}\mathbf{u}
\right),
\label{eq:proposed-regularizer}
\end{equation}
where $\Psi_{\tau}:\mathbb{R}^{|\mathcal{E}|}\rightarrow\mathbb{R}$ is defined by
\begin{equation}
\Psi_{\tau}(\mathbf{z})
:=
\inf_{\mathbf{z}=\mathbf{p}+\mathbf{q}}
\left\{
\|\mathbf{p}\|_2^2
+
2\tau\|\mathbf{q}\|_1
\right\}.
\label{eq:edge-domain-penalty}
\end{equation}
\end{definition}

In this construction, $\mathbf{p}$ is assigned a quadratic cost as in GLR, whereas $\mathbf{q}$ is assigned a linear cost, whose growth is controlled by the parameter $\tau$.
Rather than evaluating $\sqrt{\mathbf{W}}\mathbf{D}\mathbf{u}$ entirely through the quadratic cost, the optimization over its decompositions can redirect portions that would otherwise incur excessive quadratic costs to $\mathbf{q}$, for which the assigned $\ell_1$ cost grows more slowly than the quadratic cost evaluated at $\mathbf{p}$.
This mechanism depends only on the magnitudes of the proxy-weighted signal differences and therefore requires no prior knowledge of the presence, types, or locations of graph inaccuracies.

The resulting edge-domain penalty admits a standard Huber characterization through its Moreau-envelope representation \cite{Selesnick2017}.
The following fact makes this relationship explicit and will also be used in the subsequent analysis.

\begin{fact}[Moreau-envelope and Huber characterization]
\label{fact:huber}
The scaled function $\Psi_\tau/(2\tau)$ is the Moreau envelope of the $\ell_1$ norm with parameter $\tau$.
Moreover, $\Psi_\tau$ is separable and can be written as
\[
\Psi_\tau(\mathbf{z})
=
2\sum_{e\in\mathcal{E}}H_\tau(z_e),
\]
where
\[
H_\tau(t)
=
\begin{cases}
\frac{1}{2}t^2, & |t|\leq\tau,\\
\tau|t|-\frac{1}{2}\tau^2, & |t|>\tau,
\end{cases}
\]
is the scalar Huber function with threshold $\tau$.
\end{fact}

Thus, the penalty induced by the infimal-convolution design is a scaled separable Huber penalty acting on the GLR-type proxy-weighted signal differences $\sqrt{\mathbf{W}}\mathbf{D}\mathbf{u}$.
By the standard properties of Moreau envelopes, $\Psi_\tau$ is finite-valued, convex, and continuous, and so is $R_{\mathrm{IC}}$ by \eqref{eq:proposed-regularizer}.

We next use the infimal-convolution representation to characterize how the proxy-weighted signal differences are optimally decomposed between the quadratic and linear components.
This characterization makes the magnitude-adaptive allocation induced by the proposed regularizer explicit.
For notational convenience, let
\begin{equation}
\mathbf{z}
:=
\sqrt{\mathbf{W}}\mathbf{D}\mathbf{u}.
\label{eq:z-definition}
\end{equation}

\begin{proposition}[Optimal decomposition]
\label{prop:optimal-decomposition}
The infimum defining $R_{\mathrm{IC}}(\mathbf{u})$ is attained uniquely at
\begin{equation}
\mathbf{q}^{\star}
=
\operatorname{soft}(\mathbf{z};\tau),
\qquad
\mathbf{p}^{\star}
=
\operatorname{clip}(\mathbf{z};\tau),
\label{eq:optimal-decomposition}
\end{equation}
where the operators are defined componentwise as
\begin{align}
[\operatorname{soft}(\mathbf{z};\tau)]_e
&:=
\operatorname{sign}(z_e)
\max\{|z_e|-\tau,0\},
\label{eq:soft-definition}
\\
[\operatorname{clip}(\mathbf{z};\tau)]_e
&:=
\operatorname{sign}(z_e)
\min\{|z_e|,\tau\},
\label{eq:clip-definition}
\end{align}
with $\operatorname{sign}(0)=0$.
\end{proposition}

\begin{IEEEproof}
Eliminating $\mathbf{p}$ from \eqref{eq:edge-domain-penalty} shows that the infimum is attained at $\mathbf{q}^{\star}=\operatorname{prox}_{\tau\|\cdot\|_1}(\mathbf{z})$, which is the well-known soft-thresholding operator given by \eqref{eq:soft-definition}.
The equality constraint then gives $\mathbf{p}^{\star}=\mathbf{z}-\mathbf{q}^{\star}=\operatorname{clip}(\mathbf{z};\tau)$.
\end{IEEEproof}

Proposition~\ref{prop:optimal-decomposition} formalizes the intended magnitude-adaptive behavior governed by the threshold $\tau$.
If $|z_e|\leq\tau$, then $q_e^{\star}=0$, and the corresponding proxy-weighted signal difference remains entirely in the quadratically evaluated component.
If $|z_e|>\tau$, then $|p_e^{\star}|=\tau$, and the threshold-exceeding portion $|z_e|-\tau$ is assigned to $q_e^{\star}$.
The optimal decomposition further reveals how the proposed regularization value differs from the original GLR cost.

\begin{proposition}[Relation to GLR]
\label{prop:relation-to-glr}
For any $\mathbf{u}\in\mathbb{R}^{|\mathcal{V}|}$, the proposed regularizer satisfies
\begin{equation}
R_{\mathrm{IC}}(\mathbf{u})
=
R_{\mathrm{GLR}}(\mathbf{u})
-
\|\mathbf{q}^{\star}\|_2^2,
\label{eq:ric-glr-relation}
\end{equation}
where $\mathbf{q}^{\star}$ is given in Proposition~\ref{prop:optimal-decomposition}.
Moreover,
\begin{equation}
R_{\mathrm{IC}}(\mathbf{u})
=
R_{\mathrm{GLR}}(\mathbf{u})
\quad\Longleftrightarrow\quad
\left\|
\sqrt{\mathbf{W}}\mathbf{D}\mathbf{u}
\right\|_{\infty}
\leq
\tau.
\label{eq:exact-glr-coincidence}
\end{equation}
\end{proposition}

\begin{IEEEproof}
Let $\mathbf{z}:=\sqrt{\mathbf{W}}\mathbf{D}\mathbf{u}$.
Notice that the Huber function in Fact~\ref{fact:huber} can equivalently be written as
\[
2H_\tau(z_e)
=
z_e^2
-
\left(
\max\{|z_e|-\tau,0\}
\right)^2.
\]
By Proposition~\ref{prop:optimal-decomposition},
$|q_e^\star|=\max\{|z_e|-\tau,0\}$.
Hence, summing the above relation over $e\in\mathcal{E}$ and using \eqref{eq:glr-edge-domain} gives \eqref{eq:ric-glr-relation}.
Equation~\eqref{eq:ric-glr-relation} gives equality with GLR if and only if $\mathbf{q}^{\star}=\mathbf{0}$.
By Proposition~\ref{prop:optimal-decomposition}, this condition is equivalent to $\|\mathbf{z}\|_\infty\leq\tau$, which, together with \eqref{eq:z-definition}, proves \eqref{eq:exact-glr-coincidence}.
\end{IEEEproof}

Proposition~\ref{prop:relation-to-glr} shows that the proposed regularizer coincides with GLR when no proxy-weighted signal difference exceeds the threshold.
When threshold-exceeding components are present, the proposed regularizer reduces the GLR cost by $\|\mathbf{q}^{\star}\|_2^2$, thereby suppressing their excessive contribution.

\begin{remark}[Sensitivity to excessive edge weights]
\label{rem:weight-sensitivity}
By Fact~\ref{fact:huber}, the contribution associated with each proxy-weighted signal difference $z_e$ grows quadratically for $|z_e|\leq\tau$ and linearly as $2\tau|z_e|-\tau^2$ for $|z_e|>\tau$.
As an illustrative example, consider an individual proxy edge $e$ and fix its signal difference $d_e:=(\mathbf{D}\mathbf{u})_e\neq0$.
Its proxy-weighted signal difference is $z_e=\sqrt{w_e}d_e$.
Once $\sqrt{w_e}|d_e|>\tau$, the corresponding contribution to the proposed regularizer is
\begin{equation}
2\tau\sqrt{w_e}|d_e|-\tau^2.
\label{eq:weight-sensitivity-edge-cost}
\end{equation}
Hence, for fixed $d_e$, the contribution grows as $\Theta(\sqrt{w_e})$ with the proxy edge weight, whereas the corresponding GLR contribution $w_ed_e^2$ grows as $\Theta(w_e)$.
This reduced growth characterizes the sensitivity of the proposed regularizer to an overestimated edge weight after the resulting proxy-weighted signal difference exceeds the threshold.
The same argument applies to a spurious edge by interpreting $w_e>0$ as the proxy weight assigned to the newly introduced edge.
\end{remark}

Finally, we establish the smoothness of the proposed regularizer.
This property characterizes its mathematical regularity and will also be useful for optimization.

\begin{proposition}[Gradient and smoothness of the proposed regularizer]
\label{prop:mathematical-properties}
The proposed regularizer is differentiable with gradient
\begin{equation}
\nabla R_{\mathrm{IC}}(\mathbf{u})
=
2
\mathbf{D}^{\top}\sqrt{\mathbf{W}}
\operatorname{clip}
\left(
\sqrt{\mathbf{W}}\mathbf{D}\mathbf{u};
\tau
\right).
\label{eq:regularizer-gradient}
\end{equation}
Moreover, $\nabla R_{\mathrm{IC}}$ is $2\lambda_{\max}(\mathbf{L})$-Lipschitz continuous with respect to the $\ell_2$ norm.
\end{proposition}

\begin{IEEEproof}
By the Moreau-envelope characterization in Fact~\ref{fact:huber} and the standard properties of the Moreau envelope \cite[Prop.~12.30]{BauschkeCombettes2017},
\begin{equation}
\nabla\Psi_{\tau}(\mathbf{z})
=
2\left(
\mathbf{z}
-
\operatorname{prox}_{\tau\|\cdot\|_1}(\mathbf{z})
\right),
\label{eq:psi-gradient}
\end{equation}
and $\nabla\Psi_{\tau}$ is $2$-Lipschitz continuous.
Let $\mathbf{A}:=\sqrt{\mathbf{W}}\mathbf{D}$.
Applying the chain rule to \eqref{eq:proposed-regularizer} gives
\begin{equation}
\nabla R_{\mathrm{IC}}(\mathbf{u})
=
\mathbf{A}^{\top}
\nabla\Psi_{\tau}
\left(
\mathbf{A}\mathbf{u}
\right),
\end{equation}
which, together with \eqref{eq:psi-gradient} and Proposition~\ref{prop:optimal-decomposition}, yields \eqref{eq:regularizer-gradient}.
For any $\mathbf{u},\mathbf{v}\in\mathbb{R}^{|\mathcal{V}|}$, the $2$-Lipschitz continuity of $\nabla\Psi_{\tau}$ gives
\begin{align}
\|\nabla R_{\mathrm{IC}}(\mathbf{u})-\nabla R_{\mathrm{IC}}(\mathbf{v})\|_2
&\leq
2
\left\|
\mathbf{A}^{\top}
\right\|_{\mathrm{op}}
\left\|
\mathbf{A}(\mathbf{u}-\mathbf{v})
\right\|_2
\nonumber
\\
&\leq
2
\left\|
\mathbf{A}
\right\|_{\mathrm{op}}^2
\|\mathbf{u}-\mathbf{v}\|_2
\nonumber
\\
&=
2\lambda_{\max}(\mathbf{L})
\|\mathbf{u}-\mathbf{v}\|_2, 
\end{align}
where the third line follows from $\|\mathbf{A}\|_{\mathrm{op}}^2=\lambda_{\max}(\mathbf{A}^{\top}\mathbf{A})$ and \eqref{eq:laplacian}.
\end{IEEEproof}

\begin{remark}[Proximability and splitting]
\label{rem:proximability}
Although $R_{\mathrm{IC}}=\Psi_{\tau}\circ(\sqrt{\mathbf{W}}\mathbf{D})$ is not proximable in general, $\Psi_{\tau}$ itself is proximable, with its proximity operator available in closed form as given later in \eqref{eq:prox_psi}.
Accordingly, optimization problems involving $R_{\mathrm{IC}}$ can be handled by suitable splitting algorithms.
A specific instance of such an algorithm is presented in Section~IV.
\end{remark}


\subsection{Edge-Domain Sensitivity Analysis and Comparison with Potential Alternatives}
\label{subsec:alternatives}

In this subsection, we characterize the robustness of the proposed regularizer through the Lipschitz continuity of its associated edge-domain penalty $\Psi_{\tau}$.
In addition, we compare it with potential alternatives and establish an optimality property of the proposed edge-domain penalty as a replacement for the quadratic one in GLR under a prescribed Lipschitz constraint.

\begin{proposition}[Lipschitz continuity and perturbation sensitivity]
\label{prop:edge-difference-sensitivity}
For all $\mathbf{z},\widetilde{\mathbf{z}}\in\mathbb{R}^{|\mathcal{E}|}$, $\Psi_{\tau}$ is Lipschitz continuous with respect to $\ell_1$ and $\ell_2$ norms, specifically:
\begin{equation}
\left|
\Psi_{\tau}(\mathbf{z})
-
\Psi_{\tau}(\widetilde{\mathbf{z}})
\right|
\leq
2\tau\|\mathbf{z}-\widetilde{\mathbf{z}}\|_1
\leq
2\tau\sqrt{|\mathcal{E}|}\|\mathbf{z}-\widetilde{\mathbf{z}}\|_2.
\label{eq:psi-lipschitz}
\end{equation}
Consequently, for any perturbation $\boldsymbol{\delta}\in\mathbb{R}^{|\mathcal{E}|}$,
\begin{equation}
\left|
\Psi_{\tau}(\mathbf{z}+\boldsymbol{\delta})
-
\Psi_{\tau}(\mathbf{z})
\right|
\leq
2\tau\|\boldsymbol{\delta}\|_1
\leq
2\tau\sqrt{|\mathcal{E}|}\|\boldsymbol{\delta}\|_2.
\label{eq:edge-difference-sensitivity}
\end{equation}
\end{proposition}

\begin{IEEEproof}
For arbitrary $\mathbf{z}$ and $\widetilde{\mathbf{z}}$, the triangle inequality gives
\begin{align}
\Psi_{\tau}(\mathbf{z})
&=
\min_{\mathbf{y}}
\left\{
\|\mathbf{y}\|_2^2
+
2\tau\|\mathbf{z}-\mathbf{y}\|_1
\right\}
\nonumber
\\
&\leq
\min_{\mathbf{y}}
\left\{
\|\mathbf{y}\|_2^2
+
2\tau\|\widetilde{\mathbf{z}}-\mathbf{y}\|_1
\right\}
+
2\tau\|\mathbf{z}-\widetilde{\mathbf{z}}\|_1
\nonumber
\\
&=
\Psi_{\tau}(\widetilde{\mathbf{z}})
+
2\tau\|\mathbf{z}-\widetilde{\mathbf{z}}\|_1,
\end{align}
where we use $\min$ instead of $\inf$ since attainment was established in Proposition~\ref{prop:optimal-decomposition}.
Interchanging $\mathbf{z}$ and $\widetilde{\mathbf{z}}$ proves the first inequality in \eqref{eq:psi-lipschitz}, while the norm equivalence $\|\mathbf{x}\|_1\leq\sqrt{|\mathcal{E}|}\|\mathbf{x}\|_2$ gives the second.
Equation~\eqref{eq:edge-difference-sensitivity} follows by setting $\widetilde{\mathbf{z}}=\mathbf{z}+\boldsymbol{\delta}$.
\end{IEEEproof}

Equation~\eqref{eq:edge-difference-sensitivity} shows that the change in the regularization value is uniformly bounded in terms of the perturbation magnitude, independently of the original proxy-weighted signal-difference vector.
In contrast, the quadratic edge-domain penalty of GLR is not globally Lipschitz continuous and therefore does not admit such a uniform perturbation bound.

A natural alternative for obtaining bounded sensitivity is graph total variation (GTV) regularization, defined as
\begin{equation}
R_{\mathrm{GTV}}(\mathbf{u})
=
\|\mathbf{W}\mathbf{D}\mathbf{u}\|_1
=
\sum_{(i,j)\in\mathcal{E}}
w_{i,j}|u_i-u_j|.
\label{eq:gtv}
\end{equation}
Let $\mathbf{r}:=\mathbf{W}\mathbf{D}\mathbf{u}$ and $G(\mathbf{r}):=\|\mathbf{r}\|_1$ denote the weighted edge-difference vector and its edge-domain penalty, respectively.
By the triangle inequality and norm equivalence,
\begin{equation}
\left|
G(\mathbf{r})
-
G(\widetilde{\mathbf{r}})
\right|
\leq
\|\mathbf{r}-\widetilde{\mathbf{r}}\|_1
\leq
\sqrt{|\mathcal{E}|}
\|\mathbf{r}-\widetilde{\mathbf{r}}\|_2.
\label{eq:gtv-lipschitz}
\end{equation}
Thus, GTV has globally bounded sensitivity to perturbations of the weighted edge-difference vector and avoids the unbounded quadratic sensitivity of GLR to large weighted signal differences.
However, its $\ell_1$ evaluation promotes sparsity of the weighted edge differences, corresponding to a piecewise-smooth signal model rather than the quadratic graph smoothness promoted by GLR, which can lead to the well-known staircasing effect of TV-type regularization \cite{Ono2015GraphTGV}.

Since GTV provides bounded sensitivity but promotes a different signal model from GLR, another natural alternative is to supplement GLR with GTV, aiming to retain the quadratic graph-smoothness prior while incorporating the bounded-sensitivity property of the $\ell_1$ term.
We therefore consider
\begin{equation}
R_{\mathrm{GLR+GTV}}(\mathbf{u})
=
\left\|
\sqrt{\mathbf{W}}\mathbf{D}\mathbf{u}
\right\|_2^2
+
2\tau
\left\|
\mathbf{W}\mathbf{D}\mathbf{u}
\right\|_1.
\label{eq:glr-gtv}
\end{equation}
However, since the quadratic GLR term remains unchanged, the resulting regularizer still has unbounded sensitivity to large proxy-weighted signal differences, and hence the excessive-penalty issue is not eliminated.
This provides an illustrative example of the difficulty of achieving robustness through Lipschitz continuity while preserving the quadratic graph-smoothness prior of GLR.

As shown in Subsection~\ref{subsec:magnitude-adaptive}, the proposed regularizer retains the GLR evaluation for ordinary proxy-weighted signal differences while adaptively moderating the penalty growth for threshold-exceeding components, thereby achieving the global Lipschitz continuity established in Proposition~\ref{prop:edge-difference-sensitivity}.
The following proposition shows that the edge-domain penalty $\Psi_{\tau}$ associated with the proposed regularizer preserves the original quadratic GLR evaluation as faithfully as possible while satisfying the prescribed Lipschitz constraint.

\begin{proposition}[Largest Lipschitz minorant of GLR]
\label{prop:largest-lipschitz-minorant}
The function $\Psi_{\tau}$ is the pointwise largest $2\tau$-Lipschitz continuous minorant of $Q(\mathbf{z})=\|\mathbf{z}\|_2^2$ with respect to the $\ell_1$ norm.
That is, for any $2\tau$-Lipschitz continuous function $g:\mathbb{R}^{|\mathcal{E}|}\rightarrow\mathbb{R}$ satisfying $g(\mathbf{z})\leq Q(\mathbf{z})$ for all $\mathbf{z}\in\mathbb{R}^{|\mathcal{E}|}$,
\begin{equation}
g(\mathbf{z})
\leq
\Psi_{\tau}(\mathbf{z})
\leq
Q(\mathbf{z}),
\qquad
\forall\,\mathbf{z}\in\mathbb{R}^{|\mathcal{E}|}.
\label{eq:minorant-relation}
\end{equation}
\end{proposition}

\begin{IEEEproof}
By Definition~\ref{def:proposed-regularizer}, choosing $\mathbf{p}=\mathbf{z}$ and $\mathbf{q}=\mathbf{0}$ gives $\Psi_{\tau}(\mathbf{z})\leq Q(\mathbf{z})$.
Moreover, Proposition~\ref{prop:edge-difference-sensitivity} shows that $\Psi_{\tau}$ is $2\tau$-Lipschitz continuous with respect to the $\ell_1$ norm.
For any admissible $g$ and any $\mathbf{y}\in\mathbb{R}^{|\mathcal{E}|}$, the Lipschitz property and the minorant condition give
\begin{align}
g(\mathbf{z})
&\leq
g(\mathbf{y})
+
2\tau\|\mathbf{z}-\mathbf{y}\|_1
\nonumber
\\
&\leq
Q(\mathbf{y})
+
2\tau\|\mathbf{z}-\mathbf{y}\|_1.
\end{align}
Taking the minimum over $\mathbf{y}$ yields
\begin{equation}
g(\mathbf{z})
\leq
\min_{\mathbf{y}}
\left\{
Q(\mathbf{y})
+
2\tau\|\mathbf{z}-\mathbf{y}\|_1
\right\}
=
\Psi_{\tau}(\mathbf{z}),
\end{equation}
where the equality follows from Definition~\ref{def:proposed-regularizer} by setting $\mathbf{p}=\mathbf{y}$ and $\mathbf{q}=\mathbf{z}-\mathbf{y}$.
This proves \eqref{eq:minorant-relation}.
\end{IEEEproof}

In the present graph-regularization setting, this result provides a precise design interpretation of the Huber form in Fact~\ref{fact:huber}: under the prescribed Lipschitz constraint, $\Psi_\tau$ preserves the quadratic GLR edge-domain penalty pointwise as much as possible.
We include GTV and GLR+GTV in the experiments to evaluate how these different regularization mechanisms affect recovery under proxy-graph inaccuracies.